\chapter{Discussion}
\label{ch:discussion}

Chapter~\ref{ch:evaluation} answered the first two research
questions. RQ1 asked: \emph{How do varying degrees of LLM support
influence novice users during ROPA document creation?} The answer
is that more LLM support did not improve the user experience. Form,
the least LLM-driven mode, was the most usable on SUS, the
lightest on R-TLX workload, about half the wall-clock time of
either LLM-driven mode, and the most often ranked first. The
LLM-driven modes earned credit only on the AI questionnaire's
content-judgment items, with Chat highest on legal-basis
confidence; when asked which mode they would actually use again,
participants returned to Form.

RQ2 asked: \emph{How do varying degrees of LLM support influence
the quality of ROPA documents produced by novice users?} The
answer is that more LLM support did not produce categorically
better documents. The reference-based NLP metrics resolved into a
coverage-versus-precision trade-off (Form leading recall, Chat
leading semantic and precision); the chapter-level metrics split
the six ROPA chapters into three zones (Ask, Form, and Chat each
leading different sections); and the LLM-as-a-judge layer surfaced
a completeness-versus-faithfulness trade-off (Form most complete
but most fabricated, Chat most faithful but least complete). On
the judge's holistic Overall metric the three modes clustered
below the rubric midpoint with no clear winner. RQ3 asks whether
self-reported confidence aligns with measured document quality,
the cross-stream tension between RQ1's confidence findings and
RQ2's quality findings, and is what this chapter takes up: the gap
itself (Section~\ref{sec:disc-gap}), the co-author principle Ask
makes visible and the implications that follow
(Section~\ref{sec:disc-ask}), prior work
(Section~\ref{sec:disc-related}), constraints
(Section~\ref{sec:disc-limits}), and a summary
(Section~\ref{sec:disc-summary}).


\section{The Confidence--Correctness Gap}
\label{sec:disc-gap}

The study began with a prediction and the data ran the other way.
Section~\ref{sec:thesis-structure} of the Introduction set the
expected ordering as Chat $\succ$ Form $\succ$ Ask; the
post-experiment ranking placed Form first, Chat second, and Ask
third. Form led SUS, led R-TLX on time pressure and frustration,
took about half the wall-clock time of either LLM-driven mode, and
was the only mode on which the Perceived AI Support questionnaire
returned a clear lead, on re-use intent. The mode that does the
least LLM work for the user is the mode the user wants the most.

The primary self-report confidence measure was item 5 of the
Perceived AI Support questionnaire (``I am confident that the AI
identified the correct legal basis''). Participants rated Chat
highest ($3.81$), Ask second ($3.74$), and Form lowest ($3.58$).
The judge's matching quality measure, Legal Correctness, inverted
the ranking: Form $3.59$, Ask $3.43$, Chat $3.15$. Users felt most
confident in the LLM's legal judgment in the mode where the judge
rated that judgment lowest. RQ3 therefore has a direct empirical
answer on the dimension the confidence measure was designed to
probe: confidence and measured quality are inverted.

The inversion extends across the other quality dimensions. Chat
led the document-level semantic metrics (BERTScore Precision and
SBERT cosine similarity) and the judge's Faithfulness and
Hallucination-inverse ratings, but trailed on coverage (lower
BERTScore Recall, lower judge Completeness). Form led the coverage
metrics (BERTScore Recall, ROUGE-2, METEOR, and judge
Completeness) and the user-experience instruments, but the judge
rated $83.8\,\%$ of Form documents as combining maximum
Completeness with moderate-or-worse Hallucination. Preference was
highest where fabrication was highest; confidence was highest
where coverage was lowest.

The two-participant case study makes the gap concrete. IAR1116's
Chat document is the case-study ceiling on both NLP and judge
readings: highest semantic similarity to the reference, highest
scenario faithfulness on the judge rubric, and the participant's
strongest output across the metric suite. AEB2451's Form document
is the case-study floor on the judge rubric: maximum Completeness,
lowest Faithfulness and Hallucination-inverse, with invented
financial data, biometric data, video surveillance, and named
technical standards absent from the scenario. The NLP pipeline did
not flag the inventions. The Retention and Purpose passages
reproduced in
Figures~\ref{fig:worked-lexical-a}--\ref{fig:worked-semantic-b}
show both participants at passage scale; the mechanism behind the
pattern is taken up in Section~\ref{sec:disc-ask}.

The gap is not between good documents and great ones. The judge's
holistic Overall scores ($2.41$ for Form, $2.49$ for Ask, $2.62$
for Chat) all sit below the midpoint of the $1$--$5$ rubric and
more than $1.4$ points short of the ``good'' anchor at $4.0$. The
spread across modes is $0.21$ points; the gap to ``good'' is
roughly seven times that. The discussion that follows compares
three routes to a moderate document, not three deployment-ready
alternatives.

The R-TLX Performance subscale showed the calibration mismatch
inside a single instrument. Participants rated their own task
success at nearly identical levels across modes; the judge's
Completeness verdict on the same documents diverged by about a
full rubric point. Participants could not tell, from inside the
interface, which mode produced the better document.

In summary, the trade-off across the three modes runs like this.
Form covered the Article~30 fields most thoroughly and was the
mode users wanted to keep, but it scored lowest on confidence in
the LLM's legal judgment and produced the most fabricated
documents. Chat stayed closest to the scenario in semantic and
faithfulness terms and earned the highest confidence, but it
covered the fewest fields and the judge rated its legal
correctness lowest. Ask sat between the two on document quality
and worst on user experience. The user-study Form-preference and
the document-quality Chat-faithfulness do not contradict each
other; they answer different questions about the same set of
documents. The sample is more chatbot-fluent than the small or
medium-sized organization staff member the tool ultimately targets
(mostly Bachelor-level in IT or Computer Science, with high
chatbot familiarity; Section~\ref{sec:eval-demographics}), so the
Form-over-Chat preference reported here is conservative for the
broader audience the Article~30 obligation falls on
(Section~\ref{sec:background-summary}).


\section{Ask and the Co-Author Principle}
\label{sec:disc-ask}

Ask underperformed sharpest on user experience: ranked third,
slowest on completion time, most frustrating on R-TLX, and most
polarized in the free-text. On document quality the picture is
moderate; Ask sits between Form and Chat on the judge's holistic
Overall inside a narrow $0.21$-point spread. The copy-paste
friction participants reported is the visible symptom, not the
cause.

Across the three modes the LLM occupies three different positions.
In Form the user owns the structure (the checkboxes that define
which Article~30 categories the document covers) and the LLM owns
the prose. In Chat the LLM owns both. In Ask the user owns both,
and the LLM is available only on request. Ask is the only mode
where the LLM is passive, and despite that availability it made
the user do more work. To use Ask well, a novice would have needed
to know what to ask, recognize a useful answer, map fragments of a
general explanation to specific Article~30 fields, and write the
scenario-specific German legal prose. That skill set has a name:
prior GDPR knowledge.

Help that requires the user to already know what they need is not
help for a novice; asking is itself a skill. The substantive
contribution of Ask is therefore what its underperformance proves:
novice users of LLM-assisted document tooling benefit from an LLM
that co-authors or populates, not from one that advises.

The mechanism becomes visible at the implementation layer. All
three modes hand their \texttt{DocumentData} object to the same
\texttt{finalropa} prompt at temperature $0.05$ and top-$p$ $0.05$
(Methodology Table~\ref{tab:ropagen-llm-params}); only the
populated contents differ. The generator renders plausible-sounding
prose for every boolean it finds set. Form sets the most booleans,
so the generator renders the most fields and invents the most
detail. Chat sets fewer (only what the conversation extracts), so
the generator renders fewer, invents less, and omits more. Ask
sets none, so the generator works from the free-text alone. The
Form co-occurrence of maximum Completeness with moderate-or-worse
Hallucination is the trace of this mechanism, not a property of
the checkbox UI; one mechanism, the most parsimonious given the
architecture, accounts for the pattern. Hallucination at this
scale cannot be addressed by changing the interaction surface
alone, because the fabrication happens after the interface has
handed off.

The two worked examples make the mechanism visible at passage
scale. On Retention and Purpose alike, AEB2451's Form passages
outscored IAR1116's Chat passages on every reported metric, even
though Chat's content was faithful and Form's was fabricated. The
Retention passage cited $6$ to $10$ year statutory retention
against the scenario's $3$-month rule; the Purpose passage cited
legal articles (\S~257~HGB, \S~147~AO, \S~239~Abs.~1~HGB,
\S~2~Abs.~1~Nr.~1~BDSG) absent from the scenario. Reference-based
metrics measure surface and semantic similarity to a target
document; they cannot tell whether the candidate is true. BLEU
illustrates the limit sharpest: IAR1116's Chat Retention passage
shares zero $4$-grams with the reference, so BLEU-4 zeros at the
$4$-gram tier independent of the brevity penalty. The judge layer
catches what reference-based metrics structurally cannot.

The document-level comparison hides a sharper pattern at the
chapter level. The chapter-level NLP metrics split the six
rendered ROPA chapters into three zones, consistent across
lexical and BERT-based families. Ask wins Chapter~1 (Controller
and Contact Details), the most factually deterministic chapter.
Form wins Chapters~2--4 (the required-field middle), where boolean
scaffolding helps. Chat wins Chapters~5--6 (Retention and
Deletion, Technical and Organizational Measures), the interpretive
chapters whose content has to be inferred from scenario specifics.

The three-zone pattern sharpens the hybrid recommendation. The
hybrid the data supports is section-aware: route deterministic-recall
chapters to a minimal explainer, required-field chapters to
checkbox coverage, and interpretive chapters to LLM extraction
from a conversation. The simpler structure-plus-conversation
hybrid below is the necessary first step; a section-aware
refinement is the natural extension.

If novice users benefit from an LLM that co-authors rather than
advises, three concrete tooling implications follow.

The first implication is a hybrid mode that combines Form's
visible checkbox UI with Chat-style LLM extraction from a
conversation. Several participants proposed it spontaneously, and
the data supports it: Form delivers the strongest perceived
productivity, Chat delivers the strongest scenario faithfulness. A
hybrid that retains the boolean scaffolding (so the user sees what
is covered) and adds LLM extraction (so the document object
reaching generation is populated by what the user described) would
inherit both strengths. Hallucination still needs to be addressed
separately; the hybrid does not solve it.

A second implication is a validation gate at submission time. The
mechanism identified above pins the fabrication on the shared
generation step, the layer the user never sees. A validation gate
acts on that layer directly: the judge data show that a single-pass
single-model judge can already discriminate hallucinated documents
from faithful ones. A gate that runs such a judge against the
submitted document and surfaces low Faithfulness or
Hallucination-inverse scores before submit would give novice users
the only signal they currently lack. It is not a replacement for
expert review.

The third implication is an education layer for novice
calibration. Self-report confidence in this study was not
well-calibrated to document quality. An education layer that shows
novice users which fields came from their input and which came
from the LLM's generation step, with flags on fields the user did
not explicitly describe, would surface the calibration mismatch in
the interface rather than leave it for the user to detect.


\section{In Light of Prior Work}
\label{sec:disc-related}

The findings sit inside a small but active literature, surveyed in
Chapter~\ref{ch:background} along four lenses.

On GDPR compliance automation, von Schwerin et al.\
\cite{schwerinSystematicComparisonOpen2024} compared open-source
backbones for ROPA category generation and argued for
human-in-the-loop necessity on the basis of model limitations.
This thesis takes the same pipeline and extends the question from
\emph{which model} to \emph{which interaction mode}, holding the
model constant. Their human-in-the-loop claim is supported here
from a direction their study could not produce: even Chat sits
below the rubric midpoint on Overall. Abualhaija et al.\
\cite{abualhaijaLLMassistedExtractionRegulatory2025} reported, in
the adjacent problem of GDPR requirement extraction, that high
BERTScore can coexist with low BLEU when meaning is captured but
wording diverges; that pattern characterizes Chat in this study
exactly.

On structured legal-document generation, Nigam et al.\
\cite{nigamStructuredLegalDocument2025} report a
wrapper-and-iterate design for Indian private legal-document
drafting and argue that orchestration matters over scale. The
hybrid recommendation in Section~\ref{sec:disc-ask} is the
user-side counterpart; holding the model constant and varying only
the interaction surface produces sharp differences. Li et al.'s
CaseGen benchmark \cite{liCaseGenBenchmarkMultiStage2025} reports
that general-domain LLMs outperform legal-specific models, all
variants still score below $6$ out of $10$, and careful staging
matters as much as legal-domain pretraining. Both resonate here:
no mode produces what the judge calls a good document.

On evaluation methodology, Su et al.\
\cite{suJuDGEBenchmarkingJudgment2025} argued that lexical metrics
alone cannot capture legal reasoning and proposed pairing them
with judgment-specific dimensions. On this sample the genuinely
orthogonal signal is the judge, not BERT-versus-lexical: the
lexical and BERT-based families track the same per-mode ordering,
and only the judge surfaces a different ordering (Chat ahead on
Faithfulness and Hallucination-inverse, Form ahead on
Completeness). The orthogonality comes from the rubric, not from
the semantic embeddings.

On user-facing legal AI, Surya et al.\
\cite{suryaAIPoweredInteractiveLegal2025} and Vayadande et al.\
\cite{vayadandeAIPoweredLegalDocumentation2024} validated their
designs through illustrative use cases without controlled novice
studies, and both flagged that absence as a limitation. This
thesis fills the gap. The confidence--correctness decoupling is
the kind of empirical finding inspection-based evaluations cannot
produce.

Background named the unfilled intersection of three concerns:
GDPR-specific structured output, novice users without legal
training, and explicit measurement of the confidence--quality gap
under varying LLM support. This study is the first to occupy that
intersection. The substantive contribution is the gap itself, the
architectural explanation for Ask's underperformance, and the
design hierarchy that follows.


\section{Limitations}
\label{sec:disc-limits}

Five constraints bound the findings, ranging from study-design
choices to properties of the tool and model beneath it.

The first is structural. The thesis is twofold, combining a user
study with an NLP-metric evaluation, and each side carries enough
independent variables (mode, instrument, metric family, case
study, judge) that no single thread reached full depth. Findings
are broad rather than deep.

A second constraint sits in the tool. ROPAgen is a research
prototype with known bugs and architecture choices not always
optimal; Section~\ref{sec:disc-ask} already attributed Ask's
underperformance to one such choice. UI friction, prompt design,
and the boolean-scaffolding-plus-final-prose split could all shift
the measured trade-offs in a production-grade implementation.

Beyond the tool, the model layer carries a caveat. The underlying
language model is a Mistral checkpoint from $2025$, used both for
the interaction-mode dialogue and for the shared final-generation
step. Newer or larger models would likely produce different and
better documents. The relative findings about which mode helps
novices most should survive a model swap; the absolute quality
figures (all modes below the rubric midpoint) are tied to this
model.

The reference document introduces a further caveat. The NLP
metrics score against a single expert-authored ROPA whose surface
form is checkbox-derived prose, and Form's lead concentrates on
coverage metrics where reference-form match helps most. Chat leads
BERTScore Precision and SBERT, so the bias does not exhaust the
pattern. The judge, which uses no reference at all, independently
reports Form's Completeness lead. Bias is plausibly present in the
coverage metrics, not the whole story.

Finally, the judge layer is itself single-pass and single-model,
without inter-rater replication or legal-expert validation. It is
reported as a supplementary signal throughout, not subjected to
inferential testing, and triangulated against the NLP metrics and
case study. It is not a substitute for the expert review a
deployed compliance tool would require.


\section{Summary}
\label{sec:disc-summary}

The hypothesis reversed. Form, the least LLM-driven mode, was
preferred. RQ3 has a direct answer: on the legal-basis dimension
the primary confidence measure asked about, confidence and
measured quality are inverted. Chat earned the highest confidence
rating while the judge rated its legal correctness lowest; Form
earned the lowest confidence rating while the judge rated its
legal correctness highest. The wider trade-off runs across
coverage and faithfulness, with no mode best by every measure. Ask
makes the co-author principle visible by stripping the LLM of its
populator role; novice users cannot bear sole authorship without
it. The thesis-level claim that follows is that LLM-assisted
document tooling for novices should position the LLM as a
co-author, not as a teacher. Visible structure with LLM-driven
population is the hybrid the data supports; a section-aware
refinement, routing each chapter to the mode best suited to its
content, is the natural extension. The Conclusion
(Chapter~\ref{ch:conclusion}) returns to RQ1, RQ2, and RQ3 in
plain English and draws the design lessons this Discussion has
framed.
